\section{Conclusion}

We presented a calibrated, uncertainty-aware approach to continuous valence and
arousal estimation that knows when it does not know. Deep evidential regression
supplies aleatoric and epistemic uncertainty in one pass, split conformal
prediction turns that uncertainty into intervals with coverage guarantees and a
principled abstention threshold, and our evaluation centres on calibration and
distribution shift rather than accuracy alone. The same calibrated bound powers
an on-device demonstrator that abstains under poor signal and never transmits
imagery. On real AffectNet the four methods are indistinguishable on accuracy and
separate sharply on calibration, which is the result the paper rests on:
trustworthiness, not accuracy, is what the choice of head buys.

Four limitations bound the claims. Conformal is reported at a single seed.
MC-dropout and deep ensembles are released but unevaluated. The shift is
corruption-based; the cross-dataset AFEW-VA comparison is unrun and would test a
different, harder kind of shift. Most importantly, we report aggregate
calibration only. Per-group calibration is the extension that matters most,
since a threshold calibrated on an unrepresentative population fails silently for
the people it does not represent, which is the very failure this work argues
against. The gated rPPG modality, released but disabled throughout, extends the
same abstention logic to video and is the natural next step.
